% !TeX program = pdflatex
% Compilar en Overleaf con pdfLaTeX. Las figuras se incluyen en la carpeta figuras/.
% Proyecto de Sistemas Inteligentes - FES Acatlán, UNAM
\documentclass[aspectratio=169,11pt]{beamer}
\usepackage[utf8]{inputenc}
\usepackage[T1]{fontenc}
\usepackage[spanish,es-noquoting]{babel}
\usepackage{lmodern}
\usepackage{amsmath,amssymb}
\usepackage{graphicx}
\usepackage{booktabs}
\usepackage{array}
\usepackage{tikz}
\usetikzlibrary{arrows.meta,positioning,shapes.geometric,calc}
\usepackage{pgfplots}
\pgfplotsset{compat=1.18}
\usepackage{xcolor}
\usepackage{listings}
\usepackage{hyperref}

\definecolor{UNAMnavy}{HTML}{14284F}
\definecolor{UNAMdeep}{HTML}{0C1730}
\definecolor{UNAMgold}{HTML}{D2B16A}
\definecolor{UNAMteal}{HTML}{16858F}
\definecolor{Pale}{HTML}{EDF3F8}
\definecolor{Muted}{HTML}{496079}
\definecolor{Good}{HTML}{16765A}
\definecolor{Warning}{HTML}{B66B17}

\usetheme{Madrid}
\usecolortheme{default}
\setbeamercolor{structure}{fg=UNAMnavy}
\setbeamercolor{frametitle}{bg=UNAMnavy,fg=white}
\setbeamercolor{title}{fg=white}
\setbeamercolor{author in head/foot}{bg=UNAMnavy,fg=white}
\setbeamercolor{date in head/foot}{bg=UNAMnavy,fg=white}
\setbeamercolor{normal text}{fg=UNAMdeep,bg=white}
\setbeamercolor{block title}{bg=UNAMnavy,fg=white}
\setbeamercolor{block body}{bg=Pale,fg=UNAMdeep}
\setbeamercolor{alerted text}{fg=Warning}
\setbeamertemplate{navigation symbols}{}
\setbeamertemplate{footline}{%
  \leavevmode\hbox{\begin{beamercolorbox}[wd=.84\paperwidth,ht=2.5ex,dp=1ex,leftskip=1.2em]{author in head/foot}%
    Sistemas Inteligentes $\vert$ Duración de viajes NYC
  \end{beamercolorbox}\begin{beamercolorbox}[wd=.16\paperwidth,ht=2.5ex,dp=1ex,center]{date in head/foot}%
    \insertframenumber / \inserttotalframenumber
  \end{beamercolorbox}}}
\setbeamersize{text margin left=9mm,text margin right=9mm}
\setlength{\parskip}{2pt}
\setlength{\tabcolsep}{4pt}
\newcommand{\clav}[1]{\textcolor{UNAMteal}{\textbf{#1}}}
\newcommand{\porq}[1]{\vspace{1.2mm}\begin{block}{¿Por qué se hizo así?}\small #1\end{block}}
\newcommand{\grafico}[2]{\IfFileExists{figuras/#1}{\includegraphics[width=\linewidth,height=#2,keepaspectratio]{figuras/#1}}{\fbox{\parbox{.80\linewidth}{\centering Imagen opcional: figuras/#1}}}}
\newcommand{\dato}[2]{\begin{block}{#1}\centering\Large\textbf{#2}\end{block}}
\lstset{basicstyle=\ttfamily\footnotesize,backgroundcolor=\color{Pale},frame=single,rulecolor=\color{Pale},showstringspaces=false,breaklines=true,columns=fullflexible}
\hypersetup{colorlinks=true,urlcolor=UNAMteal,linkcolor=UNAMnavy}

\title[Predicción de viajes NYC]{Predicción de la duración de viajes\\de taxi en Nueva York}
\subtitle{Limpieza, ingeniería de características y selección por RMSLE}
\author{Christian Gabriel Villafuerte Aguilar}
\institute{Sistemas Inteligentes \\ Licenciatura en Matemáticas Aplicadas y Computación \\ FES Acatlán -- UNAM}
\date{Octubre de 2026}

\begin{document}

\begin{frame}[plain]
\begin{tikzpicture}[remember picture,overlay]
  \fill[UNAMnavy] (current page.south west) rectangle (current page.north east);
  \fill[UNAMgold] ([yshift=-8mm]current page.north west) rectangle ([yshift=-10mm]current page.north east);
  \node[align=center,text=white,font=\LARGE\bfseries,text width=.90\paperwidth]
      at ([yshift=12mm]current page.center)
      {Predicción de la duración\\de viajes de taxi en Nueva York};
  \node[align=center,text=UNAMgold,font=\large,text width=.90\paperwidth]
      at ([yshift=-6mm]current page.center)
      {Del CSV a un modelo de regresión evaluado por RMSLE};
  \draw[UNAMteal,line width=1.8pt]
      ([xshift=19mm,yshift=24mm]current page.south west)--([xshift=-19mm,yshift=24mm]current page.south east);
  \node[align=center,text=white,font=\normalsize]
      at ([yshift=17mm]current page.south){Christian Gabriel Villafuerte Aguilar};
  \node[align=center,text=white,font=\small]
      at ([yshift=10mm]current page.south){Sistemas Inteligentes $\cdot$ MAC $\cdot$ FES Acatlán, UNAM};
  \node[align=center,text=white,font=\small]
      at ([yshift=4mm]current page.south){Octubre de 2026};
\end{tikzpicture}
\end{frame}

\section{Comprensión del problema}
\begin{frame}{Ruta de la exposición}
\begin{enumerate}
  \item Qué se quiere predecir y cuál es la métrica oficial.
  \item Auditoría de datos, limpieza y partición correcta.
  \item Construcción de \clav{81 variables} e imputación/escalado.
  \item Comparación y ajuste de modelos de regresión.
  \item Evaluación independiente, interpretación y entrega del CSV.
\end{enumerate}
\porq{El objetivo no es mostrar solamente un puntaje: se debe justificar un proceso reproducible y evitar fuga de información.}
\end{frame}

\begin{frame}{¿Cuál es el problema?}
\begin{columns}[T]
  \column{.57\textwidth}
    \begin{itemize}
      \item Entrada: datos de cada viaje de taxi amarillo.
      \item Salida deseada: \texttt{trip\_duration}, en \textbf{segundos}.
      \item Entrenamiento: \textbf{700,000} viajes con duración conocida.
      \item Evaluación ciega: \textbf{100,000} viajes sin duración.
      \item Entrega: un valor no negativo por cada \texttt{id}.
    \end{itemize}
  \column{.39\textwidth}
    \dato{Tipo de aprendizaje}{Supervisado}
    \dato{Tipo de tarea}{Regresión}
\end{columns}
\porq{La duración puede tomar muchos valores numéricos; no estamos decidiendo una clase ``corto/largo''. Por eso no corresponde usar un clasificador.}
\end{frame}

\begin{frame}{¿Por qué RMSLE y no accuracy o F1?}
\begin{block}{Métrica oficial: Root Mean Squared Logarithmic Error}
\[
\mathrm{RMSLE}=\sqrt{\frac1n\sum_{i=1}^n\Big[\ln(1+\widehat y_i)-\ln(1+y_i)\Big]^2}.
\]
\end{block}
\begin{columns}[T]
\column{.48\textwidth}
  \textbf{En regresión:} RMSLE, MAE, RMSE, $R^2$, MAPE.
\column{.48\textwidth}
  \textbf{En clasificación:} accuracy, precision, recall, F1.
\end{columns}
\vspace{2mm}
\begin{itemize}
  \item \clav{Menor RMSLE significa mejor resultado.}
  \item La escala logarítmica da relevancia a los errores relativos.
\end{itemize}
\porq{Crear clases artificiales para calcular F1 cambiaría el problema y no optimizaría la evaluación del concurso.}
\end{frame}

\begin{frame}{Metodología de clase: ciclo de vida CRISP-DM}
\centering
\begin{tikzpicture}[>=Latex,node distance=3mm,every node/.style={font=\small}]
\tikzset{etapa/.style={draw=UNAMnavy,thick,rounded corners=3pt,fill=Pale,minimum height=11mm,text width=22mm,align=center}}
\node[etapa] (a) {Entender el\\ problema};
\node[etapa,right=of a] (b) {Auditar\\ datos};
\node[etapa,right=of b] (c) {Preparar\\ variables};
\node[etapa,right=of c] (d) {Entrenar\\ modelos};
\node[etapa,right=of d] (e) {Evaluar\\ resultados};
\draw[->,very thick,UNAMteal] (a)--(b);
\draw[->,very thick,UNAMteal] (b)--(c);
\draw[->,very thick,UNAMteal] (c)--(d);
\draw[->,very thick,UNAMteal] (d)--(e);
\draw[->,UNAMgold,thick] (e.south) .. controls +(0,-13mm) and +(0,-13mm) .. node[below,font=\scriptsize]{revisar variables y repetir si hace falta} (c.south);
\end{tikzpicture}
\vspace{8mm}
\begin{block}{Aplicación concreta en el notebook}
CSV $\rightarrow$ auditoría $\rightarrow$ partición $\rightarrow$ 81 atributos $\rightarrow$ modelos $\rightarrow$ holdout $\rightarrow$ \texttt{submission.csv}.
\end{block}
\porq{Es el enfoque explicado en la presentación \emph{Ciclo de vida del ML}: antes del algoritmo van la calidad de datos y el planteamiento de la métrica.}
\end{frame}

\section{Datos y limpieza}
\begin{frame}{¿Qué información contienen los archivos?}
\small
\begin{tabular}{p{.28\linewidth}p{.66\linewidth}}
\toprule
\textbf{Campo} & \textbf{Interpretación y uso}\tabularnewline\midrule
\texttt{id} & Identificador: se conserva para entregar, no se entrena con él.\\
\texttt{vendor\_id} & Empresa: se mantiene como variable numérica/código.\\
\texttt{pickup\_datetime} & Fecha y hora de inicio: genera variables temporales.\\
\texttt{passenger\_count} & Número declarado de pasajeros: contiene alguna anomalía.\\
\texttt{pickup\_*}, \texttt{dropoff\_*} & Coordenadas de origen/destino: permiten calcular distancias.\\
\texttt{store\_and\_fwd\_flag} & Conexión interrumpida (Y/N): se codifica 1/0.\\
\texttt{trip\_duration} & Variable objetivo, solo en \texttt{train.csv}.\\
\bottomrule
\end{tabular}
\porq{No se usa \texttt{id} como predictor porque es una etiqueta administrativa sin una relación física justificada con la duración.}
\end{frame}

\begin{frame}{Auditoría inicial: ¿los CSV están limpios?}
\begin{columns}[T]
\column{.5\textwidth}
  \begin{itemize}
    \item \textbf{0} valores nulos originales en \texttt{train} y \texttt{test}.
    \item \textbf{0} IDs duplicados en ambos archivos.
    \item \textbf{0} duraciones menores o iguales a cero.
    \item \textbf{24} registros con cero pasajeros.
  \end{itemize}
\column{.46\textwidth}
  \begin{itemize}
    \item \textbf{984} duraciones mayores que 12 horas.
    \item \textbf{2} duraciones mayores que 24 horas.
    \item \textbf{32} orígenes fuera de caja GPS NYC.
    \item \textbf{54} destinos fuera de esa caja.
  \end{itemize}
\end{columns}
\porq{Que no haya nulos no significa que todos los registros sean plausibles. También se revisan consistencia temporal, GPS y valores extremos.}
\end{frame}

\begin{frame}{Exploración de datos: antes de modelar}
\centering\grafico{01_exploracion.png}{.66\textheight}
\vspace{1mm}
{\footnotesize Distribución logarítmica de duración, número de viajes por hora y ubicaciones de origen.}
\porq{La duración tiene una cola larga; una gráfica en segundos ocultaría gran parte de los viajes usuales. Las horas y ubicaciones justifican crear variables de contexto.}
\end{frame}

\begin{frame}{Decisiones explícitas sobre anomalías}
\small
\begin{tabular}{p{.25\linewidth}p{.40\linewidth}p{.28\linewidth}}
\toprule
\textbf{Caso} & \textbf{Acción aplicada} & \textbf{Motivo}\tabularnewline\midrule
Viaje de más de 12 h & Se excluye \textbf{solo del ajuste}; no de validación ni holdout. & Reducir influencia de etiquetas extremas sin maquillar la evaluación.\\
GPS fuera de zona & Se marca con indicador y se convierte su coordenada a NaN. & No calcular distancias absurdas.\\
Cero pasajeros & Se marca con bandera de calidad; se conserva el dato. & Señal de calidad sin borrar el viaje.\\
\texttt{test.csv} & \textbf{Nunca} se eliminan filas. & Deben existir 100,000 predicciones.\\
\bottomrule
\end{tabular}
\porq{El umbral de 12 horas es una decisión heurística del experimento; conviene validarlo en estudios posteriores, no afirmar que es una regla universal.}
\end{frame}

\begin{frame}{División correcta: entrenamiento, validación y holdout}
\centering
\begin{tikzpicture}
\fill[UNAMnavy] (0,0) rectangle (10.3,1.15);
\fill[UNAMteal] (10.3,0) rectangle (11.6,1.15);
\fill[UNAMgold] (11.6,0) rectangle (12.9,1.15);
\node[white,font=\normalsize] at (5.1,.55){80\% $\rightarrow$ 560,000 viajes};
\node[white,font=\scriptsize,align=center] at (10.95,.55){10\%\\70,000};
\node[UNAMdeep,font=\scriptsize,align=center] at (12.25,.55){10\%\\70,000};
\node[below] at (5.1,-.12){Entrenamiento};
% Ambos grupos están identificados dentro de los bloques.
\end{tikzpicture}
\vspace{8mm}
\begin{itemize}
\item En los 560,000 de la partición de entrenamiento se filtraron 788 viajes $>12$ h: \clav{559,212 usados para ajustar}.
\item Validación: compara modelos y parámetros.
\item Holdout: se consulta una vez, tras la selección.
\end{itemize}
\porq{Si el holdout participa en el ajuste, deja de ser una evaluación independiente. Semilla fija: \texttt{random\_state=42}.}
\end{frame}

\section{Preparación y variables}
\begin{frame}{Del dato crudo a una matriz numérica}
\centering
\begin{tikzpicture}[>=Latex,node distance=5mm,every node/.style={font=\scriptsize}]
\tikzset{paso/.style={draw=UNAMnavy,fill=Pale,rounded corners,minimum width=22mm,minimum height=14mm,align=center}}
\node[paso] (a) {CSV\\ original};
\node[paso,right=of a] (b) {Fecha y GPS\\ consistentes};
\node[paso,right=of b] (c) {81\ atributos\\ derivados};
\node[paso,right=of c] (d) {Imputar\\ medianas};
\node[paso,right=of d] (e) {Matriz\\ $X$};
\foreach \x/\y in {a/b,b/c,c/d,d/e}{\draw[->,thick,UNAMteal] (\x)--(\y);}
\end{tikzpicture}
\vspace{8mm}
\begin{block}{Propiedad importante}
La función \texttt{create\_features} usa únicamente datos disponibles en los archivos de origen, \textbf{nunca la duración real} para construir predictores.
\end{block}
\porq{El modelo no comprende de forma automática que dos coordenadas implican una distancia o que las 23:00 y las 00:00 están muy próximas.}
\end{frame}

\begin{frame}{Las 81 características: ¿de dónde salen?}
\footnotesize
\begin{tabular}{p{.25\linewidth}p{.45\linewidth}p{.22\linewidth}}
\toprule
\textbf{Grupo} & \textbf{Ejemplos implementados} & \textbf{Qué aportan}\tabularnewline\midrule
Tiempo & Hora, minuto, día, fin de semana, hora pico, seno/coseno & Congestión y periodicidad\\
Geometría & Haversine, Manhattan aprox., $\Delta x$, $\Delta y$, rumbo & Longitud y dirección\\
Zonas & Rejillas GPS, puntos medios & Barrios y trayectos frecuentes\\
Referencias & Distancia a JFK, LGA, EWR, Times Square y otros & Contexto urbano\\
Interacciones & Hora $\times$ distancia, hora pico $\times$ distancia & Relaciones no lineales\\
Calidad & GPS inválido, pasajero inválido, conexión diferida & Anomalías del registro\\
\bottomrule
\end{tabular}
\porq{Estas 81 columnas resultan de combinar conocimiento del dominio con transformaciones reproducibles; su cantidad no garantiza por sí sola calidad.}
\end{frame}

\begin{frame}{Variables temporales: ¿por qué usar seno y coseno?}
\begin{columns}[T]
\column{.52\textwidth}
\[
 h_{\sin}=\sin\left(\frac{2\pi h}{24}\right),\qquad
 h_{\cos}=\cos\left(\frac{2\pi h}{24}\right).
\]
\begin{itemize}
\item 23:00 y 00:00 son horas contiguas.
\item Una hora numérica normal no refleja bien esa cercanía.
\item También se aplican ciclos a días y meses.
\end{itemize}
\column{.44\textwidth}
\centering
\begin{tikzpicture}[scale=1.15]
\draw[->,gray] (-1.35,0)--(1.35,0);
\draw[->,gray] (0,-1.35)--(0,1.35);
\draw[thick,UNAMnavy] (0,0) circle (1);
\fill[UNAMteal] (1,0) circle (.07);
\fill[UNAMgold] ({cos(-15)},{sin(-15)}) circle (.075);
\node[font=\scriptsize,right] at (1,.12){00 h};
\node[font=\scriptsize,right] at (0.84,-.37){23 h};
\end{tikzpicture}
\end{columns}
\porq{La codificación cíclica hace que el cambio de día no aparezca como un salto artificial de 23 a 0.}
\end{frame}

\begin{frame}{¿Cómo calcular la distancia entre dos coordenadas?}
\begin{block}{Fórmula de Haversine (distancia sobre la esfera terrestre)}
\small
\[
\begin{aligned}
 a&=\sin^2\!\left(\frac{\Delta\varphi}{2}\right)
 +\cos\varphi_1\cos\varphi_2\sin^2\!\left(\frac{\Delta\lambda}{2}\right),\\[-1mm]
 d&=2R\arcsin(\sqrt{a}), \qquad R\simeq 6371\,\mathrm{km}.
\end{aligned}
\]
\end{block}
\begin{itemize}
\item $\varphi$: latitud; $\lambda$: longitud, en radianes.
\item La distancia Haversine estima el trayecto directo, \textbf{no} las calles recorridas.
\item Se añade distancia Manhattan aproximada para representar desplazamientos ortogonales.
\end{itemize}
\porq{Las coordenadas originales son difíciles de interpretar solas; las distancias aportan una relación más directa con el tiempo esperado.}
\end{frame}

\begin{frame}{Geografía urbana y relaciones entre variables}
\begin{columns}[T]
\column{.49\textwidth}
\begin{block}{Origen / destino}
\begin{itemize}
\item Distancia a JFK, LaGuardia y Newark.
\item Distancia a Times Square, Grand Central y Downtown.
\item Celdas y punto medio de la trayectoria.
\end{itemize}
\end{block}
\column{.49\textwidth}
\begin{block}{Combinaciones}
\begin{itemize}
\item \texttt{hour\_x\_distance}.
\item \texttt{rush\_x\_distance}.
\item Ángulo del trayecto y sus seno/coseno.
\item Distancias al cuadrado y logarítmicas.
\end{itemize}
\end{block}
\end{columns}
\porq{Un trayecto de 5 km durante hora pico puede durar distinto que el mismo trayecto por la madrugada; la interacción captura esa dependencia.}
\end{frame}

\begin{frame}{Imputación: ¿qué hacer cuando una coordenada es inválida?}
\begin{enumerate}
\item Detectar coordenadas fuera de una caja geográfica razonable.
\item Crear bandera de calidad y sustituir la coordenada inválida por \texttt{NaN}.
\item Aprender la \textbf{mediana de cada característica} usando \textbf{únicamente entrenamiento}.
\item Aplicar las mismas medianas a validación, holdout y test.
\end{enumerate}
\begin{block}{Implementación real}
\texttt{SimpleImputer(strategy='median', keep\_empty\_features=True)}
\end{block}
\porq{La mediana suele ser más resistente que la media a extremos. Si se calculase usando todo el conjunto, entrarían estadísticas de validación o prueba en el entrenamiento.}
\end{frame}

\begin{frame}{Normalización, escalado y logaritmo: NO son lo mismo}
\small
\begin{tabular}{p{.28\linewidth}p{.32\linewidth}p{.30\linewidth}}
\toprule
\textbf{Técnica} & \textbf{Expresión} & \textbf{¿Se usó?}\tabularnewline\midrule
Estandarización & $(x-\mu)/\sigma$ & Se explica, pero no se aplicó a los árboles.\\
Min--max & $(x-x_{\min})/(x_{\max}-x_{\min})$ & No se aplicó.\\
Escalado robusto & $(x-\mathrm{mediana})/\mathrm{IQR}$ & \textbf{Sí: solo en Ridge}.\\
Transformación del objetivo & $z=\ln(1+y)$ & \textbf{Sí: en todos los modelos}.\\
\bottomrule
\end{tabular}
\vspace{2mm}
\porq{Ridge penaliza coeficientes y necesita escalas comparables; los árboles separan por umbrales, así que no requieren escalado para funcionar.}
\end{frame}

\begin{frame}{Transformación de la duración para entrenar}
\begin{columns}[T]
\column{.52\textwidth}
\begin{block}{Objetivo de entrenamiento}
\[
 z=\ln(1+\texttt{trip\_duration}).
\]
\end{block}
\begin{block}{Recuperar segundos}
\[
 \widehat y=\max\{0,\exp(\widehat z)-1\}.
\]
\end{block}
\column{.44\textwidth}
\begin{itemize}
\item Reduce la asimetría de la duración.
\item Es consistente con la métrica RMSLE.
\item Evita valores negativos en el archivo final.
\item \textbf{No} significa que se entregue el logaritmo.
\end{itemize}
\end{columns}
\porq{Se optimiza en la escala logarítmica y después se regresa a segundos, que es la unidad requerida por el reto.}
\end{frame}

\begin{frame}{¿Cómo se evitó la fuga de información?}
\centering
\begin{tikzpicture}[node distance=7mm,>=Latex,every node/.style={font=\small}]
\node[draw=UNAMnavy,thick,fill=Pale,rounded corners,minimum width=41mm,minimum height=10mm] (tr) {Entrenamiento \texttt{tr\_fit}};
\node[draw=UNAMteal,thick,fill=Pale,rounded corners,right=17mm of tr,minimum width=53mm,minimum height=10mm] (imp) {\texttt{imputer.fit(Xtr)}};
\draw[->,thick,UNAMteal] (tr)--(imp);
\node[draw=UNAMnavy,fill=white,below=5mm of tr,minimum width=41mm,minimum height=8mm] (va) {Validación / holdout / test};
\node[draw=UNAMteal,fill=white,below=5mm of imp,minimum width=53mm,minimum height=8mm] (tx) {\texttt{imputer.transform(X)}};
\draw[->,thick,UNAMteal] (va)--(tx);
\draw[->,thick,UNAMgold] (imp)--(tx);
\end{tikzpicture}
\vspace{8mm}
\begin{itemize}
\item No se ajustan medianas ni \texttt{RobustScaler} con validación o test.
\item No se calcula velocidad usando duración real ni se inventa hora de llegada.
\item La evaluación holdout ocurre después de seleccionar el ganador.
\end{itemize}
\porq{Una métrica artificialmente baja por fuga de información no mide la capacidad real de generalizar.}
\end{frame}

\section{Modelado y selección}
\begin{frame}{Primero: dos referencias fáciles de explicar}
\begin{columns}[T]
\column{.48\textwidth}
\begin{block}{Modelo ingenuo: mediana}
Predice para todos los viajes la mediana de la duración en escala logarítmica.
\vspace{2mm}
\centering\Large\textbf{RMSLE = 0.795413}
\end{block}
\column{.48\textwidth}
\begin{block}{Ridge: regresión lineal}
Modelo lineal con regularización $L_2$ y \texttt{RobustScaler}.
\vspace{2mm}
\centering\Large\textbf{RMSLE = 0.462847}
\end{block}
\end{columns}
\porq{El baseline fija el mínimo que debemos superar; Ridge evalúa cuánto se puede explicar con una relación lineal regularizada.}
\end{frame}

\begin{frame}{¿Por qué Ridge en vez de quedarse con regresión lineal?}
\begin{block}{Idea matemática}
\[
\widehat{\boldsymbol\beta}
 =\arg\min_{\boldsymbol\beta}\left\{\sum_{i=1}^n(z_i-\beta_0-\mathbf x_i^\top\boldsymbol\beta)^2
 +\alpha\lVert\boldsymbol\beta\rVert_2^2\right\},\qquad \alpha=100.
\]
\end{block}
\begin{itemize}
\item La penalización reduce coeficientes excesivamente grandes.
\item Es útil con muchas variables correlacionadas.
\item Se ajusta sobre atributos escalados de forma robusta.
\item Pero tiene dificultades con interacciones no lineales complejas.
\end{itemize}
\porq{Ridge es una referencia interpretable y regularizada, no necesariamente el modelo final.}
\end{frame}

\begin{frame}{¿Por qué pasar a árboles de boosting?}
\begin{columns}[T]
\column{.51\textwidth}
\begin{itemize}
\item Distancia y duración no mantienen siempre relación lineal.
\item Influyen hora, zona, aeropuerto y trayecto simultáneamente.
\item \textbf{Boosting} añade árboles que corrigen errores anteriores.
\item Puede capturar interacciones y cortes no lineales.
\end{itemize}
\column{.45\textwidth}
\centering
\begin{tikzpicture}[scale=.87,>=Latex]
\node[draw=UNAMnavy,fill=Pale,rounded corners] (a) at (0,1.3){Árbol 1};
\node[draw=UNAMnavy,fill=Pale,rounded corners] (b) at (2.2,1.3){Árbol 2};
\node[draw=UNAMnavy,fill=Pale,rounded corners] (c) at (4.4,1.3){Árbol 3};
\draw[->,thick,UNAMteal] (a)--(b);
\draw[->,thick,UNAMteal] (b)--(c);
\draw[->,thick,UNAMgold] (2.2,.8)--(2.2,0) node[below,font=\small,align=center]{suma gradual\\de correcciones};
\end{tikzpicture}
\end{columns}
\porq{Los métodos de árboles trabajan con umbrales; por ello aquí no añadimos escalado de entrada que no aporta una ventaja clara.}
\end{frame}

\begin{frame}{Modelos probados y motivo de comparación}
\footnotesize
\begin{tabular}{p{.23\linewidth}p{.43\linewidth}p{.27\linewidth}}
\toprule
\textbf{Modelo} & \textbf{Idea} & \textbf{Motivo del ensayo}\tabularnewline\midrule
LightGBM & Boosting eficiente con árboles orientados a hojas. & Eficiencia en datos tabulares grandes.\\
XGBoost & Boosting regularizado con árboles. & Comparación con otra implementación sólida.\\
CatBoost & Boosting con mecanismos de regularización propios. & Contrastar estabilidad y desempeño.\\
\midrule
Blends & Combinaciones ponderadas en escala log. & Comprobar si se complementan sus errores.\\
\bottomrule
\end{tabular}
\porq{No se declaró ganador por popularidad: todos se compararon con el mismo conjunto de validación y la misma métrica.}
\end{frame}

\begin{frame}{¿Qué parámetros se ajustaron en LightGBM?}
\scriptsize
\begin{tabular}{p{.36\linewidth}rrr}
\toprule
\textbf{Hiperparámetro} & \textbf{Equilibrado} & \textbf{Profundo} & \textbf{Regularizado}\\\midrule
\texttt{n\_estimators} & 1400 & 1600 & 1650\\
\texttt{learning\_rate} & 0.045 & 0.045 & 0.035\\
\texttt{num\_leaves} & 63 & 105 & 39\\
\texttt{min\_child\_samples} & 45 & 75 & 75\\
\texttt{colsample\_bytree} & 0.87 & 0.92 & 0.83\\
\texttt{subsample} & 0.87 & 0.90 & 0.88\\
\texttt{reg\_lambda} & 2.5 & 5 & 6\\\midrule
\textbf{RMSLE de validación} & \textbf{0.379249} & \textbf{0.375735} & \textbf{0.382771}\\
\bottomrule
\end{tabular}
\vspace{2mm}
\begin{block}{Criterio aplicado}
Entrenamiento sobre $\ln(1+y)$; evaluación por RMSLE; \emph{early stopping} con 90 rondas sin mejora para LightGBM.
\end{block}
\porq{Más hojas pueden capturar estructuras complejas, pero se controlan con mínimo de muestras por hoja y regularización. No se exploró todo el espacio de parámetros.}
\end{frame}

\begin{frame}{Resultados de validación: comparación objetiva}
\begin{columns}[T]
\column{.53\textwidth}
\scriptsize
\begin{tabular}{lr}
\toprule Modelo & RMSLE $\downarrow$\\\midrule
\textbf{LightGBM profundo} & \textbf{0.375735}\\
LightGBM equilibrado & 0.379249\\
XGBoost & 0.381371\\
LightGBM regularizado & 0.382771\\
CatBoost & 0.385040\\
Ridge escalado & 0.462847\\
Mediana & 0.795413\\
\bottomrule
\end{tabular}
\par\vspace{3mm}
\noindent\textbf{Ganador entre los modelos ensayados:} LightGBM profundo.
\column{.43\textwidth}
\begin{tikzpicture}
\begin{axis}[xbar,xmin=0,xmax=.85,width=\linewidth,height=54mm,bar width=7pt,axis x line=bottom,axis y line=none,ytick=data,yticklabels={Mediana,Ridge,CatBoost,XGBoost,LightGBM},tick label style={font=\scriptsize},nodes near coords,point meta=explicit symbolic,every node near coord/.append style={font=\tiny,anchor=west},enlarge y limits=.14]
\addplot[fill=UNAMteal,draw=UNAMteal] coordinates {(.7954,0) [.795] (.4628,1) [.463] (.3850,2) [.385] (.3814,3) [.381] (.3757,4) [.376]};
\end{axis}
\end{tikzpicture}
\end{columns}
\porq{El único criterio de selección fue minimizar RMSLE en la validación; valores menores indican mejores predicciones relativas.}
\end{frame}

\begin{frame}{¿Se intentó combinar modelos? Sí}
\small
\begin{tabular}{p{.75\linewidth}r}
\toprule
\textbf{Ensamble probado} & \textbf{RMSLE}\\\midrule
75\% LightGBM profundo + 25\% LightGBM equilibrado & 0.375912\\
50\% LightGBM profundo + 50\% LightGBM equilibrado & 0.376559\\
25\% LightGBM profundo + 75\% LightGBM equilibrado & 0.377673\\
50/30/20 (profundo / equilibrado / XGBoost) & 0.376475\\
\midrule
\textbf{LightGBM profundo sin mezcla} & \textbf{0.375735}\\\bottomrule
\end{tabular}
\vspace{2mm}
\porq{Un ensamble no garantiza mejora. Como ninguno superó al modelo individual, el resultado final usa un solo LightGBM: mejor RMSLE y menor complejidad.}
\end{frame}

\section{Evaluación y entrega}
\begin{frame}{Evaluación independiente: 70,000 viajes holdout}
\begin{columns}[T]
\column{.49\textwidth}
\dato{RMSLE en validación}{0.375735}
\column{.49\textwidth}
\dato{RMSLE en holdout}{0.381219}
\end{columns}
\vspace{2mm}
\begin{center}
\small
\begin{tabular}{ll}
\toprule\textbf{Métrica complementaria en holdout} & \textbf{Valor}\\\midrule
MAE & 270.92 segundos\\
RMSE & 3,062.55 segundos\\
$R^2$ & 0.0390\\
MAPE & 25.48\%\\\bottomrule
\end{tabular}
\end{center}
\porq{El holdout se mantuvo ajeno a la selección, por lo que su RMSLE es una estimación más honesta de generalización; no es el puntaje de test ciego.}
\end{frame}

\begin{frame}{Interpretación: RMSLE vs. MAE, RMSE y $R^2$}
\small
\begin{tabular}{p{.19\linewidth}p{.71\linewidth}}
\toprule
\textbf{Métrica} & \textbf{Qué nos dice}\tabularnewline\midrule
RMSLE & Penaliza errores al comparar \textbf{log-duraciones}; es la métrica oficial.\\
MAE & En promedio hay \textbf{270.92 s} de error absoluto en holdout (unos 4.5 min).\\
RMSE & \textbf{3,062.55 s}; muy sensible a errores grandes en segundos.\\
$R^2$ & \textbf{0.039}: el ajuste en escala de segundos explica poca varianza de estas etiquetas.\\
MAPE & \textbf{25.48\%}, pero puede ser inestable cuando la duración real es pequeña.\\
\bottomrule
\end{tabular}
\porq{El mismo modelo puede rendir bien en RMSLE y mal en RMSE/$R^2$: no son métricas intercambiables y los valores extremos influyen mucho en segundos.}
\end{frame}

\begin{frame}{Diagnóstico: ¿en qué casos se equivoca?}
\centering\grafico{03_diagnostico_holdout.png}{.55\textheight}
\vspace{1mm}
{\footnotesize Predicción frente a valor real en logaritmos; distribución del error logarítmico.}
\porq{Las dispersiones y sesgos ayudan a detectar sobreestimación, subestimación y viajes atípicos. Una sola métrica agregada no explica todos los fallos.}
\end{frame}

\begin{frame}{Importancia de variables según LightGBM}
\begin{columns}[T]
\column{.54\textwidth}
\centering\grafico{04_importancia_variables.png}{.61\textheight}
\column{.43\textwidth}
\small
\begin{itemize}
\item Distancia a \textbf{JFK} desde el origen.
\item Componentes cíclicas de la \textbf{hora}.
\item \textbf{Día del año}.
\item Distancia a JFK desde destino.
\item \textbf{Hora $\times$ distancia}.
\item \textbf{Distancia Haversine}.
\end{itemize}
\end{columns}
\porq{Es una importancia interna del árbol (cantidad de uso de variables en cortes), no demuestra causalidad ni ordena necesariamente el efecto real de cada característica.}
\end{frame}

\begin{frame}{El error depende de la distancia del trayecto}
\small
\begin{tabular}{lrr}
\toprule\textbf{Distancia Haversine} & \textbf{Viajes holdout} & \textbf{RMSLE}\\\midrule
Menos de 0.5 km & 2,504 & 1.103559\\
0.5 a 2 km & 31,070 & 0.368468\\
2 a 5 km & 23,896 & 0.290010\\
5 a 10 km & 8,250 & 0.260631\\
10 a 25 km & 4,200 & 0.275809\\
Más de 25 km & 76 & 0.245508\\\bottomrule
\end{tabular}
\vspace{1mm}
{\scriptsize Nota: cuatro viajes holdout no entran en estos rangos por GPS inválido.}\par
\vspace{1mm}
\porq{En viajes muy cortos, paradas, espera y registros de GPS pueden dominar el tiempo; algunos grupos tienen pocas observaciones. Se trata de un diagnóstico, no de una prueba causal.}
\end{frame}

\begin{frame}{Entrenamiento final, ahora sí con todos los datos utilizables}
\begin{enumerate}
\item Se seleccionó LightGBM profundo usando \textbf{solo validación}.
\item Se verificó su desempeño en \textbf{holdout}.
\item Se volvió a entrenar con \textbf{699,016 viajes etiquetados}, aplicando el filtro $1\leq y\leq 43,200$ s.
\item Se calcularon de nuevo medianas con ese entrenamiento completo.
\item Se guardó el modelo: \texttt{modelos/modelo\_final.joblib}.
\end{enumerate}
\porq{Después de cerrar la evaluación, tiene sentido aprovechar más ejemplos etiquetados. El modelo final ya no tiene un holdout propio medido con esos pesos nuevos.}
\end{frame}

\begin{frame}[fragile]{Generación y verificación del archivo de entrega}
\begin{columns}[T]
\column{.49\textwidth}
\textbf{Formato exigido:}
\begin{lstlisting}
id,trip_duration
id_00001,650.79453
id_00002,412.00000
...
\end{lstlisting}
\footnotesize Los números anteriores son \textbf{ilustrativos del formato}, no filas copiadas de la entrega.
\column{.47\textwidth}
\textbf{Comprobaciones programadas:}
\begin{itemize}\small
\item \textbf{100,000} registros.
\item Columnas exactas y en orden.
\item Cada \texttt{id} aparece una sola vez.
\item Mismo orden que \texttt{test.csv}.
\item Segundos numéricos, finitos y $\geq 0$.
\end{itemize}
\end{columns}
\porq{No basta con entrenar bien: un CSV con IDs omitidos, columnas de más o valores negativos sería una entrega incorrecta.}
\end{frame}

\begin{frame}{Conclusiones principales}
\begin{enumerate}
\item El problema es \textbf{regresión}, por lo que el objetivo correcto es \textbf{minimizar RMSLE}.
\item La auditoría detectó duraciones extremas y coordenadas GPS incoherentes.
\item La creación de \textbf{81 variables} permite traducir contexto temporal y geográfico a números.
\item \textbf{Ridge se escaló}; los modelos de árbol \textbf{no necesitaron escalado}.
\item LightGBM profundo ganó entre los candidatos: \clav{0.375735 en validación} y \clav{0.381219 en holdout}.
\item El archivo \texttt{submission.csv} tiene las 100,000 predicciones requeridas.
\end{enumerate}
\end{frame}

\begin{frame}{Limitaciones y trabajo futuro}
\begin{itemize}
\item No conocemos la duración real del \textbf{test ciego}: no podemos asegurar su RMSLE.
\item El \textbf{RMSLE de holdout} es una estimación, no una garantía de puntaje.
\item El umbral de \textbf{12 horas} es una decisión a contrastar con otras estrategias.
\item No se dispone aquí de tráfico real, accidentes, clima o rutas de calles.
\item Podrían probarse \textbf{validación temporal}, optimización más amplia de hiperparámetros y variables de tránsito; habría que volver a evaluar sin fuga.
\end{itemize}
\porq{``Mejor modelo'' significa \textbf{el mejor entre los ensayados}, no una demostración de óptimo global.}
\end{frame}

\begin{frame}{Preguntas que pueden hacer en clase}
\small
\begin{enumerate}
\item \textbf{¿Por qué no F1 o accuracy?} Porque se predicen segundos continuos.
\item \textbf{¿Por qué el logaritmo?} Por cola larga y por coherencia con RMSLE.
\item \textbf{¿Por qué no escalaste LightGBM?} Porque las particiones por umbrales no dependen de esa escala.
\item \textbf{¿Por qué filtrar $>12$ h solo en train?} Para limitar ruido sin ocultarlo en evaluación.
\item \textbf{¿Por qué LightGBM y no XGBoost?} Tuvo menor RMSLE de validación en estos experimentos.
\item \textbf{¿Conoces el puntaje de test?} No; carece de etiquetas de duración.
\end{enumerate}
\end{frame}

\begin{frame}{Fuentes y archivos para reproducir}
\small
\begin{itemize}
\item \emph{U5 - 01 - ML.pdf}: aprendizaje supervisado; regresión vs. clasificación (p. 7).
\item \emph{U5 - 02 - Perceptrón.pdf}: regresión lineal, función de costo, descenso de gradiente (pp. 3--7).
\item \emph{03 - Ciclo de vida del ML.pdf}: CRISP-DM, calidad, partición, fuga, imputación, escalado y métricas (pp. 2--5, 8--15, 19--24).
\item \texttt{Prediccion\_Taxis\_NYC\_Paso\_a\_Paso.ipynb}: ejecución y decisiones exactas.
\item \texttt{resultados/metricas.json}, \texttt{comparacion\_modelos.csv} y gráficas: evidencia experimental.
\end{itemize}
\vspace{2mm}
\begin{center}\Large\textbf{¡Gracias!}\end{center}
\end{frame}

\end{document}
